\section{谱测度把“对角化”推广到连续能量}

先看已经对角化的无穷维算子。\(\mathcal H=L^2(\mathbb R)\) 上令 \((Mf)(x)=xf(x)\)，定义域为 \(D(M)=\{f\in L^2:xf\in L^2\}\)。这个域稠密，因为紧支光滑函数都在其中。它也恰好等于伴随域：若 \(g\in D(M^*)\)，按伴随定义存在 \(h\in L^2\) 使 \(\int\overline g,xf=\int\overline h f\) 对所有 \(f\in D(M)\) 成立；先让 \(f\) 遍历紧支光滑函数，就得到 \(xg=h\) 几乎处处。因此 \(xg\in L^2\)，即 \(g\in D(M)\)；反向由乘法规则直接成立。故 \(M=M^*\)，不是只在形式上对称。

它却没有本征向量：若 \(Mf=\lambda f\)，则 \((x-\lambda)f(x)=0\) 几乎处处，故 \(f\) 只能支撑在单点；单点测度为零，所以 \(f=0\) 作为 \(L^2\) 向量。按谱的定义，所有实数仍在 \(\sigma(M)\)。确实，对实 \(\lambda\) 取支撑在 \((\lambda-1/n,\lambda+1/n)\) 且范数为一的 \(f_n\)，有 \(\|(M-\lambda)f_n\|\le1/n\)。若 \((M-\lambda)^{-1}\) 有界，则 \(1\le\|(M-\lambda)^{-1}\|/n\)，矛盾。反之 \(z\notin\mathbb R\) 时 \((M-z)^{-1}f(x)=f(x)/(x-z)\) 存在且范数至多 \(|\operatorname{Im}z|^{-1}\)，所以 \(\sigma(M)=\mathbb R\)。这里每个谱点都是“近似本征值”，却没有任何真正的 \(L^2\) 本征向量；这正是有限维直觉遗漏的情形。

在这个例子中，\((E(B)f)(x)=\mathbf1_B(x)f(x)\)。它确实是投影，因为 \(\mathbf1_B^2=\mathbf1_B\)，并且 \(E(B)E(C)=E(B\cap C)\)。对 \(f\) 而言，\(\mu_f(B)=\langle f,E(B)f\rangle=\int_B|f(x)|^2\,\dd x\) 是一个非负测度。于是
\[
\|Mf\|^2=\int_{\mathbb R}\lambda^2\,\dd\mu_f(\lambda),
\qquad (h(M)f)(x)=h(x)f(x).
\]
这里 \(\mu_f\) 完全由普通积分给出，不需要虚构一个“本征函数 \(\delta(x-\lambda)\)”作为 Hilbert 空间向量。

\begin{theorem}[自伴算子的谱定理]
每个自伴算子 \(A\) 都有唯一的投影值测度 \(E_A\)，使 \(A=\int\lambda\,\dd E_A(\lambda)\)。对 Borel 函数 \(h:\mathbb R\to\mathbb C\)，定义
\[
D(h(A))=\left\{x:\int|h(\lambda)|^2\,\dd\mu_x(\lambda)<\infty\right\},
\qquad \mu_x(B)=\langle x,E_A(B)x\rangle,
\]
且 \(\|h(A)x\|^2=\int|h|^2\,\dd\mu_x\)。若 \(h\) 有界，则 \(h(A)\) 在整个 \(\mathcal H\) 上有界。
\end{theorem}

一旦知道 \(E_A\)，怎样从投影真的构造 \(h(A)\)？对有限个互不相交能量集合上的阶梯函数 \(h=\sum_{j=1}^N a_j\mathbf1_{B_j}\)，先定义 \(h(A)x=\sum_j a_jE_A(B_j)x\)。投影分量正交，故
\[
 \|h(A)x\|^2=\sum_j|a_j|^2\mu_x(B_j)
 =\int|h|^2\,\dd\mu_x.
\]

对一般 Borel 函数 \(h\)，取一列与 \(x\) 无关、由截断及复平面网格化构成的阶梯函数 \(h_N\)，使 \(h_N(\lambda)\to h(\lambda)\)，且 \(|h_N(\lambda)|\le|h(\lambda)|+1\)。若 \(\int|h|^2\,\dd\mu_x<\infty\)，\((|h|+1)^2\) 对有限测度 \(\mu_x\) 可积，支配收敛定理给 \(\int|h_N-h_M|^2\,\dd\mu_x\to0\)。于是 \(h_N(A)x\) 为 Cauchy 列；其极限定义 \(h(A)x\)。范数恒等式还保证换用另一种阶梯逼近不会改变极限。谱定理最深的部分是保证每个自伴 \(A\) 都有唯一的 \(E_A\)；上面的构造则说明为什么有了它便能一致地计算演化、预解式和能量概率。乘法算子的计算已经逐项验证了存在部分在具体函数空间中的每个对象。

取 \(h(\lambda)=\ee^{-\ii t\lambda/\hbar}\)，其模恒为一，所以 \(U(t)=h(A)\) 对所有态有定义且保持范数。取 \(h(\lambda)=(\lambda-z)^{-1}\)，若 \(\operatorname{Im}z\ne0\)，则 \(|h|\le1/|\operatorname{Im}z|\)，从而
\[
(A-zI)^{-1}=\int\frac{1}{\lambda-z}\,\dd E_A(\lambda),
\qquad \|(A-zI)^{-1}\|\le\frac1{|\operatorname{Im}z|}.
\]
同一测度同时给出演化与求源响应，原因只是给谱变量选择了不同函数。

\section{从预解式读回谱的哪一部分}

对乘法算子，\(R(z)f(x)=f(x)/(x-z)\)。当 \(z\) 靠近实轴时，分母小的区域贡献最强。但“靠近”究竟能恢复哪一块能量概率，必须先把极限次序说清。对任意固定 \(\varepsilon>0\)，谱表示和 \(\mu_f(\mathbb R)=\|f\|^2<\infty\) 给出
\[
 \operatorname{Im}\langle f,(A-\lambda-\ii\varepsilon)^{-1}f\rangle
 =\int_{\mathbb R}\frac{\varepsilon}{(t-\lambda)^2+\varepsilon^2}\,\dd\mu_f(t).
\]
被积函数非负，所以可交换 \(\lambda\) 与 \(t\) 的积分。对固定 \(t\)，内层积分是
\[
 P_\varepsilon(t;a,b)=\frac1\pi\int_a^b
 \frac{\varepsilon}{(t-\lambda)^2+\varepsilon^2}\,\dd\lambda
 =\frac1\pi\left[\arctan\frac{b-t}{\varepsilon}
 -\arctan\frac{a-t}{\varepsilon}\right].
\]
它始终处于 \([0,1]\)，而且 \(\varepsilon\downarrow0\) 时在 \(a<t<b\) 趋于一，在区间外趋于零，在 \(t=a,b\) 趋于二分之一。因为 \(\mu_f\) 是有限测度，有界收敛定理允许把极限送进 \(t\) 积分。因此严格的反演式是
\begin{equation}
 \lim_{\varepsilon\downarrow0}\frac1\pi\int_a^b
 \operatorname{Im}\langle f,(A-\lambda-\ii\varepsilon)^{-1}f\rangle\,\dd\lambda
 =\mu_f((a,b))+\frac12\mu_f(\{a\})+\frac12\mu_f(\{b\}).
 \label{eq:math-stieltjes-inversion}
\end{equation}
若两端都没有原子，它才简化为 \(\mu_f((a,b))\)。这个论证同时覆盖离散谱、连续谱及二者混合的情形，不要求 \(\mu_f\) 有密度。孤立本征值 \(\lambda_0\) 的权重可由包含它而不含其他谱点的小区间取出，即 \(\mu_f(\{\lambda_0\})=\|E_A(\{\lambda_0\})f\|^2\)；若在区间端点取样却忘记半权重，就会差一倍。

只有在 \(\mu_f\) 对 Lebesgue 测度绝对连续、写得成 \(\dd\mu_f(t)=\rho_f(t)\,\dd t\) 时，才谈得上普通意义的谱密度 \(\rho_f\)。在 \(\rho_f\) 的连续点，Poisson 核的近似恒等性质进一步给 \(\rho_f(\lambda)=\pi^{-1}\lim_{\varepsilon\downarrow0}\operatorname{Im}\langle f,(A-\lambda-\ii\varepsilon)^{-1}f\rangle\)。对原子谱，同一个虚部在本征值处随 \(\varepsilon^{-1}\) 发散，不能把它误读为有限的函数值；先对能量区间积分才是两类谱共同适用的做法。

这个极限可在已经求出的区间 Laplacian 上逐项检查。令 \(S_D\phi_n=n^2\pi^2\phi_n\)，\(\phi_n=\sqrt2\sin(n\pi x)\)，并把 \(f=\sum c_n\phi_n\)。则
\[
E_{S_D}(B)f=\sum_{n^2\pi^2\in B}c_n\phi_n,
\qquad
\langle f,(S_D-z)^{-1}f\rangle
=\sum_{n\ge1}\frac{|c_n|^2}{n^2\pi^2-z}.
\]
若 \((a,b)\) 只含 \(m^2\pi^2\)，把后一式的虚部对 \(\lambda\in(a,b)\) 积分；第 \(m\) 项的 Poisson 核趋于 \(\pi|c_m|^2\)，其余项趋零，于是恢复 \(\|E_{S_D}((a,b))f\|^2=|c_m|^2\)。相反，乘法算子的 \(\mu_f(B)=\int_B|f|^2\) 没有孤立原子。投影值测度在两种谱中是同一种对象，差别只在 \(\mu_f\) 是原子和还是连续积分。

预解式与 Green 核也在这里接上：对 Dirichlet Laplacian 的 \(z=0\)，\((S_D-0)^{-1}f(x)=\int_0^1G(x,y)f(y)\,\dd y\)。把正弦展开代入左侧，便得到 \(G(x,y)=\sum_{n\ge1}\phi_n(x)\phi_n(y)/(n^2\pi^2)\) 的核表示；它与下一节由端点齐次解拼出的 \(\min(x,y)[1-\max(x,y)]\) 是同一函数。前一个表达看见谱，后一个表达看见边界与导数跳跃。

\begin{exercise}
对 \(M f(x)=xf(x)\)，取 \(f=\mathbf1_{[0,1]}\)。直接求 \(\mu_f((a,b))\) 与 \(\|\ee^{-\ii tM}f\|\)，再验证 \(R(z)f=f/(x-z)\)。
\end{exercise}
